% =============================================================================
%  Fundamentos compartilhados --- Módulo 01: ferramentas de desenvolvimento
%  Incluído pelas trilhas Kotlin e Flutter com % =============================================================================
%  Fundamentos compartilhados --- Módulo 01: ferramentas de desenvolvimento
%  Incluído pelas trilhas Kotlin e Flutter com % =============================================================================
%  Fundamentos compartilhados --- Módulo 01: ferramentas de desenvolvimento
%  Incluído pelas trilhas Kotlin e Flutter com % =============================================================================
%  Fundamentos compartilhados --- Módulo 01: ferramentas de desenvolvimento
%  Incluído pelas trilhas Kotlin e Flutter com \input{../fundamentos/f01-ambiente}
% =============================================================================
\section{Fundamentos: ferramentas de desenvolvimento}

\begin{frame}{Do código-fonte ao programa}
  \centering
  \begin{tikzpicture}[node distance=5mm,every node/.style={font=\scriptsize,align=center,minimum height=10mm}]
    \node[caixa,text width=1.9cm] (src) {Código-fonte\\ \tiny texto legível};
    \node[caixapetroleo,right=of src,text width=1.9cm] (comp) {Compilador\\ \tiny verifica e traduz};
    \node[caixaambar,right=of comp,text width=2.2cm] (bc) {\emph{Bytecode} ou\\ código de máquina};
    \node[caixapetroleo,right=of bc,text width=1.9cm] (pk) {Empacotador\\ \tiny junta recursos};
    \node[caixaescura,right=of pk,text width=1.9cm] (app) {App\\ instalável};
    \foreach \a/\b in {src/comp,comp/bc,bc/pk,pk/app} \draw[seta] (\a) -- (\b);
  \end{tikzpicture}
  \vspace{4pt}
  \begin{columns}[T]
    \begin{column}{0.48\textwidth}
      \begin{definicao}[Compilador e interpretador]
        O \textbf{compilador} traduz o programa inteiro antes da execução e aponta erros cedo. O \textbf{interpretador} lê e executa o código aos poucos, durante a execução.
      \end{definicao}
      \begin{definicao}[Máquina virtual]
        Programa que executa um \emph{bytecode} intermediário em qualquer processador: a JVM e o ART (Kotlin) e a VM do Dart.
      \end{definicao}
    \end{column}
    \begin{column}{0.48\textwidth}
      \begin{definicao}[JIT e AOT]
        \textbf{JIT}: compila enquanto o programa roda; permite trocar código em execução (\emph{hot reload}).
        \textbf{AOT}: compila tudo antes de distribuir; inicia rápido e roda com desempenho previsível.
      \end{definicao}
      \begin{dica}
        As duas trilhas usam os dois: JIT para desenvolver com agilidade, AOT para o app que chega ao usuário (Módulo 15).
      \end{dica}
    \end{column}
  \end{columns}
\end{frame}

\begin{frame}{SDK, IDE e cadeia de ferramentas}
  \begin{columns}[T]
    \begin{column}{0.48\textwidth}
      \begin{definicao}[SDK --- \emph{Software Development Kit}]
        Pacote com bibliotecas, compiladores, ferramentas de linha de comando, documentação e imagens de emulador para desenvolver para uma plataforma.
      \end{definicao}
      \begin{definicao}[IDE --- ambiente integrado]
        Editor que reúne escrita de código, compilação, execução, depuração, análise e integração com Git.
      \end{definicao}
      \begin{analogia}
        O SDK é a caixa de ferramentas; a IDE é a bancada organizada onde você trabalha com elas.
      \end{analogia}
    \end{column}
    \begin{column}{0.5\textwidth}
      \scriptsize
      \begin{tabularx}{\linewidth}{@{}l L L@{}}
        \toprule
        & \cabecalho{Trilha Kotlin} & \cabecalho{Trilha Flutter} \\
        \midrule
        Linguagem & Kotlin (sobre o JDK) & Dart (Dart SDK) \\
        SDK da plataforma & Android SDK & Flutter SDK + Android SDK + Xcode \\
        IDE & Android Studio & VS Code ou Android Studio \\
        \emph{Build} & Gradle & \cod{flutter} CLI (usa Gradle e Xcode) \\
        Aparelho & ADB & ADB e \cod{flutter devices} \\
        Pacotes & Maven Central, Google Maven & pub.dev \\
        \bottomrule
      \end{tabularx}
    \end{column}
  \end{columns}
\end{frame}

\begin{frame}{Sistemas de \emph{build} e dependências}
  \begin{columns}[c]
    \begin{column}{0.46\textwidth}
      \centering
      \begin{tikzpicture}[node distance=4mm,every node/.style={font=\tiny,minimum width=1.6cm}]
        \node[caixaescura] (apk) {gerar pacote};
        \node[caixapetroleo,below left=5mm and -2mm of apk] (comp) {compilar};
        \node[caixapetroleo,below right=5mm and -2mm of apk] (res) {processar recursos};
        \node[caixa,below=5mm of comp] (dep) {baixar dependências};
        \node[caixa,below=5mm of res] (gen) {gerar código};
        \draw[seta] (comp) -- (apk); \draw[seta] (res) -- (apk);
        \draw[seta] (dep) -- (comp); \draw[seta] (gen) -- (comp);
      \end{tikzpicture}\\[2pt]
      {\tiny Grafo de tarefas: cada uma só roda depois das que ela usa.}
    \end{column}
    \begin{column}{0.52\textwidth}
      \begin{definicao}[Sistema de \emph{build}]
        Ferramenta que transforma o projeto em app executando \textbf{tarefas} organizadas num grafo de dependências. Só refaz o que mudou (\emph{build} incremental) e guarda resultados em cache.
      \end{definicao}
      \begin{definicao}[Dependência transitiva]
        A biblioteca que você usa depende de outras, que dependem de outras. O gerenciador resolve a árvore inteira e escolhe versões compatíveis.
      \end{definicao}
      \begin{atencao}
        Fixe as versões com arquivos de trava (catálogo de versões, \cod{pubspec.lock}) para que todos na equipe compilem o mesmo app.
      \end{atencao}
    \end{column}
  \end{columns}
\end{frame}

\begin{frame}{Emuladores, simuladores e aparelhos reais}
  \begin{columns}[T]
    \begin{column}{0.5\textwidth}
      \scriptsize
      \begin{tabularx}{\linewidth}{@{}l L L@{}}
        \toprule
        & \cabecalho{Vantagem} & \cabecalho{Limite} \\
        \midrule
        Emulador Android & vários tamanhos e versões, GPS e sensores simulados & sem Bluetooth, NFC e USB reais \\
        Simulador iOS & rápido, integrado ao Xcode & roda código do Mac; não imita o hardware \\
        Aparelho real & desempenho e hardware verdadeiros & cada modelo exige um aparelho \\
        \bottomrule
      \end{tabularx}
    \end{column}
    \begin{column}{0.48\textwidth}
      \begin{definicao}[Emulador e simulador]
        O \textbf{emulador} imita o aparelho, inclusive o processador e o sistema completo. O \textbf{simulador} apenas reproduz o comportamento do sistema, executando o app compilado para o computador.
      \end{definicao}
      \begin{pordentro}[Por que o emulador exige virtualização]
        O emulador Android roda uma imagem completa do sistema numa máquina virtual. Com aceleração por hardware (Intel VT-x, AMD-V, KVM no Linux, Hypervisor no macOS), as instruções rodam direto no processador e o emulador fica rápido.
      \end{pordentro}
    \end{column}
  \end{columns}
\end{frame}

\begin{frame}{Depuração}
  \begin{columns}[T]
    \begin{column}{0.5\textwidth}
      \begin{definicao}[Depurador]
        Ferramenta que pausa o programa num ponto escolhido (\emph{breakpoint}) para inspecionar variáveis, a pilha de chamadas e executar passo a passo.
      \end{definicao}
      \scriptsize
      \begin{tabularx}{\linewidth}{@{}l L@{}}
        \toprule
        \cabecalho{Comando} & \cabecalho{O que faz} \\
        \midrule
        \emph{Step over} & executa a linha e para na próxima \\
        \emph{Step into} & entra na função chamada \\
        \emph{Step out} & termina a função atual e volta \\
        \emph{Resume} & continua até o próximo \emph{breakpoint} \\
        \bottomrule
      \end{tabularx}
    \end{column}
    \begin{column}{0.48\textwidth}
      \begin{conceito}[Método científico da depuração]
        \begin{enumerate}
          \item Reproduza o problema de forma confiável.
          \item Formule uma hipótese sobre a causa.
          \item Teste a hipótese (\emph{breakpoint}, \emph{log}, teste automatizado).
          \item Corrija e escreva um teste que impeça a volta do defeito.
        \end{enumerate}
      \end{conceito}
      \begin{dica}
        Leia a mensagem de erro e a pilha de chamadas inteiras: a primeira linha do \textbf{seu} código costuma apontar a causa.
      \end{dica}
    \end{column}
  \end{columns}
\end{frame}

% =============================================================================
\section{Fundamentos: ferramentas de desenvolvimento}

\begin{frame}{Do código-fonte ao programa}
  \centering
  \begin{tikzpicture}[node distance=5mm,every node/.style={font=\scriptsize,align=center,minimum height=10mm}]
    \node[caixa,text width=1.9cm] (src) {Código-fonte\\ \tiny texto legível};
    \node[caixapetroleo,right=of src,text width=1.9cm] (comp) {Compilador\\ \tiny verifica e traduz};
    \node[caixaambar,right=of comp,text width=2.2cm] (bc) {\emph{Bytecode} ou\\ código de máquina};
    \node[caixapetroleo,right=of bc,text width=1.9cm] (pk) {Empacotador\\ \tiny junta recursos};
    \node[caixaescura,right=of pk,text width=1.9cm] (app) {App\\ instalável};
    \foreach \a/\b in {src/comp,comp/bc,bc/pk,pk/app} \draw[seta] (\a) -- (\b);
  \end{tikzpicture}
  \vspace{4pt}
  \begin{columns}[T]
    \begin{column}{0.48\textwidth}
      \begin{definicao}[Compilador e interpretador]
        O \textbf{compilador} traduz o programa inteiro antes da execução e aponta erros cedo. O \textbf{interpretador} lê e executa o código aos poucos, durante a execução.
      \end{definicao}
      \begin{definicao}[Máquina virtual]
        Programa que executa um \emph{bytecode} intermediário em qualquer processador: a JVM e o ART (Kotlin) e a VM do Dart.
      \end{definicao}
    \end{column}
    \begin{column}{0.48\textwidth}
      \begin{definicao}[JIT e AOT]
        \textbf{JIT}: compila enquanto o programa roda; permite trocar código em execução (\emph{hot reload}).
        \textbf{AOT}: compila tudo antes de distribuir; inicia rápido e roda com desempenho previsível.
      \end{definicao}
      \begin{dica}
        As duas trilhas usam os dois: JIT para desenvolver com agilidade, AOT para o app que chega ao usuário (Módulo 15).
      \end{dica}
    \end{column}
  \end{columns}
\end{frame}

\begin{frame}{SDK, IDE e cadeia de ferramentas}
  \begin{columns}[T]
    \begin{column}{0.48\textwidth}
      \begin{definicao}[SDK --- \emph{Software Development Kit}]
        Pacote com bibliotecas, compiladores, ferramentas de linha de comando, documentação e imagens de emulador para desenvolver para uma plataforma.
      \end{definicao}
      \begin{definicao}[IDE --- ambiente integrado]
        Editor que reúne escrita de código, compilação, execução, depuração, análise e integração com Git.
      \end{definicao}
      \begin{analogia}
        O SDK é a caixa de ferramentas; a IDE é a bancada organizada onde você trabalha com elas.
      \end{analogia}
    \end{column}
    \begin{column}{0.5\textwidth}
      \scriptsize
      \begin{tabularx}{\linewidth}{@{}l L L@{}}
        \toprule
        & \cabecalho{Trilha Kotlin} & \cabecalho{Trilha Flutter} \\
        \midrule
        Linguagem & Kotlin (sobre o JDK) & Dart (Dart SDK) \\
        SDK da plataforma & Android SDK & Flutter SDK + Android SDK + Xcode \\
        IDE & Android Studio & VS Code ou Android Studio \\
        \emph{Build} & Gradle & \cod{flutter} CLI (usa Gradle e Xcode) \\
        Aparelho & ADB & ADB e \cod{flutter devices} \\
        Pacotes & Maven Central, Google Maven & pub.dev \\
        \bottomrule
      \end{tabularx}
    \end{column}
  \end{columns}
\end{frame}

\begin{frame}{Sistemas de \emph{build} e dependências}
  \begin{columns}[c]
    \begin{column}{0.46\textwidth}
      \centering
      \begin{tikzpicture}[node distance=4mm,every node/.style={font=\tiny,minimum width=1.6cm}]
        \node[caixaescura] (apk) {gerar pacote};
        \node[caixapetroleo,below left=5mm and -2mm of apk] (comp) {compilar};
        \node[caixapetroleo,below right=5mm and -2mm of apk] (res) {processar recursos};
        \node[caixa,below=5mm of comp] (dep) {baixar dependências};
        \node[caixa,below=5mm of res] (gen) {gerar código};
        \draw[seta] (comp) -- (apk); \draw[seta] (res) -- (apk);
        \draw[seta] (dep) -- (comp); \draw[seta] (gen) -- (comp);
      \end{tikzpicture}\\[2pt]
      {\tiny Grafo de tarefas: cada uma só roda depois das que ela usa.}
    \end{column}
    \begin{column}{0.52\textwidth}
      \begin{definicao}[Sistema de \emph{build}]
        Ferramenta que transforma o projeto em app executando \textbf{tarefas} organizadas num grafo de dependências. Só refaz o que mudou (\emph{build} incremental) e guarda resultados em cache.
      \end{definicao}
      \begin{definicao}[Dependência transitiva]
        A biblioteca que você usa depende de outras, que dependem de outras. O gerenciador resolve a árvore inteira e escolhe versões compatíveis.
      \end{definicao}
      \begin{atencao}
        Fixe as versões com arquivos de trava (catálogo de versões, \cod{pubspec.lock}) para que todos na equipe compilem o mesmo app.
      \end{atencao}
    \end{column}
  \end{columns}
\end{frame}

\begin{frame}{Emuladores, simuladores e aparelhos reais}
  \begin{columns}[T]
    \begin{column}{0.5\textwidth}
      \scriptsize
      \begin{tabularx}{\linewidth}{@{}l L L@{}}
        \toprule
        & \cabecalho{Vantagem} & \cabecalho{Limite} \\
        \midrule
        Emulador Android & vários tamanhos e versões, GPS e sensores simulados & sem Bluetooth, NFC e USB reais \\
        Simulador iOS & rápido, integrado ao Xcode & roda código do Mac; não imita o hardware \\
        Aparelho real & desempenho e hardware verdadeiros & cada modelo exige um aparelho \\
        \bottomrule
      \end{tabularx}
    \end{column}
    \begin{column}{0.48\textwidth}
      \begin{definicao}[Emulador e simulador]
        O \textbf{emulador} imita o aparelho, inclusive o processador e o sistema completo. O \textbf{simulador} apenas reproduz o comportamento do sistema, executando o app compilado para o computador.
      \end{definicao}
      \begin{pordentro}[Por que o emulador exige virtualização]
        O emulador Android roda uma imagem completa do sistema numa máquina virtual. Com aceleração por hardware (Intel VT-x, AMD-V, KVM no Linux, Hypervisor no macOS), as instruções rodam direto no processador e o emulador fica rápido.
      \end{pordentro}
    \end{column}
  \end{columns}
\end{frame}

\begin{frame}{Depuração}
  \begin{columns}[T]
    \begin{column}{0.5\textwidth}
      \begin{definicao}[Depurador]
        Ferramenta que pausa o programa num ponto escolhido (\emph{breakpoint}) para inspecionar variáveis, a pilha de chamadas e executar passo a passo.
      \end{definicao}
      \scriptsize
      \begin{tabularx}{\linewidth}{@{}l L@{}}
        \toprule
        \cabecalho{Comando} & \cabecalho{O que faz} \\
        \midrule
        \emph{Step over} & executa a linha e para na próxima \\
        \emph{Step into} & entra na função chamada \\
        \emph{Step out} & termina a função atual e volta \\
        \emph{Resume} & continua até o próximo \emph{breakpoint} \\
        \bottomrule
      \end{tabularx}
    \end{column}
    \begin{column}{0.48\textwidth}
      \begin{conceito}[Método científico da depuração]
        \begin{enumerate}
          \item Reproduza o problema de forma confiável.
          \item Formule uma hipótese sobre a causa.
          \item Teste a hipótese (\emph{breakpoint}, \emph{log}, teste automatizado).
          \item Corrija e escreva um teste que impeça a volta do defeito.
        \end{enumerate}
      \end{conceito}
      \begin{dica}
        Leia a mensagem de erro e a pilha de chamadas inteiras: a primeira linha do \textbf{seu} código costuma apontar a causa.
      \end{dica}
    \end{column}
  \end{columns}
\end{frame}

% =============================================================================
\section{Fundamentos: ferramentas de desenvolvimento}

\begin{frame}{Do código-fonte ao programa}
  \centering
  \begin{tikzpicture}[node distance=5mm,every node/.style={font=\scriptsize,align=center,minimum height=10mm}]
    \node[caixa,text width=1.9cm] (src) {Código-fonte\\ \tiny texto legível};
    \node[caixapetroleo,right=of src,text width=1.9cm] (comp) {Compilador\\ \tiny verifica e traduz};
    \node[caixaambar,right=of comp,text width=2.2cm] (bc) {\emph{Bytecode} ou\\ código de máquina};
    \node[caixapetroleo,right=of bc,text width=1.9cm] (pk) {Empacotador\\ \tiny junta recursos};
    \node[caixaescura,right=of pk,text width=1.9cm] (app) {App\\ instalável};
    \foreach \a/\b in {src/comp,comp/bc,bc/pk,pk/app} \draw[seta] (\a) -- (\b);
  \end{tikzpicture}
  \vspace{4pt}
  \begin{columns}[T]
    \begin{column}{0.48\textwidth}
      \begin{definicao}[Compilador e interpretador]
        O \textbf{compilador} traduz o programa inteiro antes da execução e aponta erros cedo. O \textbf{interpretador} lê e executa o código aos poucos, durante a execução.
      \end{definicao}
      \begin{definicao}[Máquina virtual]
        Programa que executa um \emph{bytecode} intermediário em qualquer processador: a JVM e o ART (Kotlin) e a VM do Dart.
      \end{definicao}
    \end{column}
    \begin{column}{0.48\textwidth}
      \begin{definicao}[JIT e AOT]
        \textbf{JIT}: compila enquanto o programa roda; permite trocar código em execução (\emph{hot reload}).
        \textbf{AOT}: compila tudo antes de distribuir; inicia rápido e roda com desempenho previsível.
      \end{definicao}
      \begin{dica}
        As duas trilhas usam os dois: JIT para desenvolver com agilidade, AOT para o app que chega ao usuário (Módulo 15).
      \end{dica}
    \end{column}
  \end{columns}
\end{frame}

\begin{frame}{SDK, IDE e cadeia de ferramentas}
  \begin{columns}[T]
    \begin{column}{0.48\textwidth}
      \begin{definicao}[SDK --- \emph{Software Development Kit}]
        Pacote com bibliotecas, compiladores, ferramentas de linha de comando, documentação e imagens de emulador para desenvolver para uma plataforma.
      \end{definicao}
      \begin{definicao}[IDE --- ambiente integrado]
        Editor que reúne escrita de código, compilação, execução, depuração, análise e integração com Git.
      \end{definicao}
      \begin{analogia}
        O SDK é a caixa de ferramentas; a IDE é a bancada organizada onde você trabalha com elas.
      \end{analogia}
    \end{column}
    \begin{column}{0.5\textwidth}
      \scriptsize
      \begin{tabularx}{\linewidth}{@{}l L L@{}}
        \toprule
        & \cabecalho{Trilha Kotlin} & \cabecalho{Trilha Flutter} \\
        \midrule
        Linguagem & Kotlin (sobre o JDK) & Dart (Dart SDK) \\
        SDK da plataforma & Android SDK & Flutter SDK + Android SDK + Xcode \\
        IDE & Android Studio & VS Code ou Android Studio \\
        \emph{Build} & Gradle & \cod{flutter} CLI (usa Gradle e Xcode) \\
        Aparelho & ADB & ADB e \cod{flutter devices} \\
        Pacotes & Maven Central, Google Maven & pub.dev \\
        \bottomrule
      \end{tabularx}
    \end{column}
  \end{columns}
\end{frame}

\begin{frame}{Sistemas de \emph{build} e dependências}
  \begin{columns}[c]
    \begin{column}{0.46\textwidth}
      \centering
      \begin{tikzpicture}[node distance=4mm,every node/.style={font=\tiny,minimum width=1.6cm}]
        \node[caixaescura] (apk) {gerar pacote};
        \node[caixapetroleo,below left=5mm and -2mm of apk] (comp) {compilar};
        \node[caixapetroleo,below right=5mm and -2mm of apk] (res) {processar recursos};
        \node[caixa,below=5mm of comp] (dep) {baixar dependências};
        \node[caixa,below=5mm of res] (gen) {gerar código};
        \draw[seta] (comp) -- (apk); \draw[seta] (res) -- (apk);
        \draw[seta] (dep) -- (comp); \draw[seta] (gen) -- (comp);
      \end{tikzpicture}\\[2pt]
      {\tiny Grafo de tarefas: cada uma só roda depois das que ela usa.}
    \end{column}
    \begin{column}{0.52\textwidth}
      \begin{definicao}[Sistema de \emph{build}]
        Ferramenta que transforma o projeto em app executando \textbf{tarefas} organizadas num grafo de dependências. Só refaz o que mudou (\emph{build} incremental) e guarda resultados em cache.
      \end{definicao}
      \begin{definicao}[Dependência transitiva]
        A biblioteca que você usa depende de outras, que dependem de outras. O gerenciador resolve a árvore inteira e escolhe versões compatíveis.
      \end{definicao}
      \begin{atencao}
        Fixe as versões com arquivos de trava (catálogo de versões, \cod{pubspec.lock}) para que todos na equipe compilem o mesmo app.
      \end{atencao}
    \end{column}
  \end{columns}
\end{frame}

\begin{frame}{Emuladores, simuladores e aparelhos reais}
  \begin{columns}[T]
    \begin{column}{0.5\textwidth}
      \scriptsize
      \begin{tabularx}{\linewidth}{@{}l L L@{}}
        \toprule
        & \cabecalho{Vantagem} & \cabecalho{Limite} \\
        \midrule
        Emulador Android & vários tamanhos e versões, GPS e sensores simulados & sem Bluetooth, NFC e USB reais \\
        Simulador iOS & rápido, integrado ao Xcode & roda código do Mac; não imita o hardware \\
        Aparelho real & desempenho e hardware verdadeiros & cada modelo exige um aparelho \\
        \bottomrule
      \end{tabularx}
    \end{column}
    \begin{column}{0.48\textwidth}
      \begin{definicao}[Emulador e simulador]
        O \textbf{emulador} imita o aparelho, inclusive o processador e o sistema completo. O \textbf{simulador} apenas reproduz o comportamento do sistema, executando o app compilado para o computador.
      \end{definicao}
      \begin{pordentro}[Por que o emulador exige virtualização]
        O emulador Android roda uma imagem completa do sistema numa máquina virtual. Com aceleração por hardware (Intel VT-x, AMD-V, KVM no Linux, Hypervisor no macOS), as instruções rodam direto no processador e o emulador fica rápido.
      \end{pordentro}
    \end{column}
  \end{columns}
\end{frame}

\begin{frame}{Depuração}
  \begin{columns}[T]
    \begin{column}{0.5\textwidth}
      \begin{definicao}[Depurador]
        Ferramenta que pausa o programa num ponto escolhido (\emph{breakpoint}) para inspecionar variáveis, a pilha de chamadas e executar passo a passo.
      \end{definicao}
      \scriptsize
      \begin{tabularx}{\linewidth}{@{}l L@{}}
        \toprule
        \cabecalho{Comando} & \cabecalho{O que faz} \\
        \midrule
        \emph{Step over} & executa a linha e para na próxima \\
        \emph{Step into} & entra na função chamada \\
        \emph{Step out} & termina a função atual e volta \\
        \emph{Resume} & continua até o próximo \emph{breakpoint} \\
        \bottomrule
      \end{tabularx}
    \end{column}
    \begin{column}{0.48\textwidth}
      \begin{conceito}[Método científico da depuração]
        \begin{enumerate}
          \item Reproduza o problema de forma confiável.
          \item Formule uma hipótese sobre a causa.
          \item Teste a hipótese (\emph{breakpoint}, \emph{log}, teste automatizado).
          \item Corrija e escreva um teste que impeça a volta do defeito.
        \end{enumerate}
      \end{conceito}
      \begin{dica}
        Leia a mensagem de erro e a pilha de chamadas inteiras: a primeira linha do \textbf{seu} código costuma apontar a causa.
      \end{dica}
    \end{column}
  \end{columns}
\end{frame}

% =============================================================================
\section{Fundamentos: ferramentas de desenvolvimento}

\begin{frame}{Do código-fonte ao programa}
  \centering
  \begin{tikzpicture}[node distance=5mm,every node/.style={font=\scriptsize,align=center,minimum height=10mm}]
    \node[caixa,text width=1.9cm] (src) {Código-fonte\\ \tiny texto legível};
    \node[caixapetroleo,right=of src,text width=1.9cm] (comp) {Compilador\\ \tiny verifica e traduz};
    \node[caixaambar,right=of comp,text width=2.2cm] (bc) {\emph{Bytecode} ou\\ código de máquina};
    \node[caixapetroleo,right=of bc,text width=1.9cm] (pk) {Empacotador\\ \tiny junta recursos};
    \node[caixaescura,right=of pk,text width=1.9cm] (app) {App\\ instalável};
    \foreach \a/\b in {src/comp,comp/bc,bc/pk,pk/app} \draw[seta] (\a) -- (\b);
  \end{tikzpicture}
  \vspace{4pt}
  \begin{columns}[T]
    \begin{column}{0.48\textwidth}
      \begin{definicao}[Compilador e interpretador]
        O \textbf{compilador} traduz o programa inteiro antes da execução e aponta erros cedo. O \textbf{interpretador} lê e executa o código aos poucos, durante a execução.
      \end{definicao}
      \begin{definicao}[Máquina virtual]
        Programa que executa um \emph{bytecode} intermediário em qualquer processador: a JVM e o ART (Kotlin) e a VM do Dart.
      \end{definicao}
    \end{column}
    \begin{column}{0.48\textwidth}
      \begin{definicao}[JIT e AOT]
        \textbf{JIT}: compila enquanto o programa roda; permite trocar código em execução (\emph{hot reload}).
        \textbf{AOT}: compila tudo antes de distribuir; inicia rápido e roda com desempenho previsível.
      \end{definicao}
      \begin{dica}
        As duas trilhas usam os dois: JIT para desenvolver com agilidade, AOT para o app que chega ao usuário (Módulo 15).
      \end{dica}
    \end{column}
  \end{columns}
\end{frame}

\begin{frame}{SDK, IDE e cadeia de ferramentas}
  \begin{columns}[T]
    \begin{column}{0.48\textwidth}
      \begin{definicao}[SDK --- \emph{Software Development Kit}]
        Pacote com bibliotecas, compiladores, ferramentas de linha de comando, documentação e imagens de emulador para desenvolver para uma plataforma.
      \end{definicao}
      \begin{definicao}[IDE --- ambiente integrado]
        Editor que reúne escrita de código, compilação, execução, depuração, análise e integração com Git.
      \end{definicao}
      \begin{analogia}
        O SDK é a caixa de ferramentas; a IDE é a bancada organizada onde você trabalha com elas.
      \end{analogia}
    \end{column}
    \begin{column}{0.5\textwidth}
      \scriptsize
      \begin{tabularx}{\linewidth}{@{}l L L@{}}
        \toprule
        & \cabecalho{Trilha Kotlin} & \cabecalho{Trilha Flutter} \\
        \midrule
        Linguagem & Kotlin (sobre o JDK) & Dart (Dart SDK) \\
        SDK da plataforma & Android SDK & Flutter SDK + Android SDK + Xcode \\
        IDE & Android Studio & VS Code ou Android Studio \\
        \emph{Build} & Gradle & \cod{flutter} CLI (usa Gradle e Xcode) \\
        Aparelho & ADB & ADB e \cod{flutter devices} \\
        Pacotes & Maven Central, Google Maven & pub.dev \\
        \bottomrule
      \end{tabularx}
    \end{column}
  \end{columns}
\end{frame}

\begin{frame}{Sistemas de \emph{build} e dependências}
  \begin{columns}[c]
    \begin{column}{0.46\textwidth}
      \centering
      \begin{tikzpicture}[node distance=4mm,every node/.style={font=\tiny,minimum width=1.6cm}]
        \node[caixaescura] (apk) {gerar pacote};
        \node[caixapetroleo,below left=5mm and -2mm of apk] (comp) {compilar};
        \node[caixapetroleo,below right=5mm and -2mm of apk] (res) {processar recursos};
        \node[caixa,below=5mm of comp] (dep) {baixar dependências};
        \node[caixa,below=5mm of res] (gen) {gerar código};
        \draw[seta] (comp) -- (apk); \draw[seta] (res) -- (apk);
        \draw[seta] (dep) -- (comp); \draw[seta] (gen) -- (comp);
      \end{tikzpicture}\\[2pt]
      {\tiny Grafo de tarefas: cada uma só roda depois das que ela usa.}
    \end{column}
    \begin{column}{0.52\textwidth}
      \begin{definicao}[Sistema de \emph{build}]
        Ferramenta que transforma o projeto em app executando \textbf{tarefas} organizadas num grafo de dependências. Só refaz o que mudou (\emph{build} incremental) e guarda resultados em cache.
      \end{definicao}
      \begin{definicao}[Dependência transitiva]
        A biblioteca que você usa depende de outras, que dependem de outras. O gerenciador resolve a árvore inteira e escolhe versões compatíveis.
      \end{definicao}
      \begin{atencao}
        Fixe as versões com arquivos de trava (catálogo de versões, \cod{pubspec.lock}) para que todos na equipe compilem o mesmo app.
      \end{atencao}
    \end{column}
  \end{columns}
\end{frame}

\begin{frame}{Emuladores, simuladores e aparelhos reais}
  \begin{columns}[T]
    \begin{column}{0.5\textwidth}
      \scriptsize
      \begin{tabularx}{\linewidth}{@{}l L L@{}}
        \toprule
        & \cabecalho{Vantagem} & \cabecalho{Limite} \\
        \midrule
        Emulador Android & vários tamanhos e versões, GPS e sensores simulados & sem Bluetooth, NFC e USB reais \\
        Simulador iOS & rápido, integrado ao Xcode & roda código do Mac; não imita o hardware \\
        Aparelho real & desempenho e hardware verdadeiros & cada modelo exige um aparelho \\
        \bottomrule
      \end{tabularx}
    \end{column}
    \begin{column}{0.48\textwidth}
      \begin{definicao}[Emulador e simulador]
        O \textbf{emulador} imita o aparelho, inclusive o processador e o sistema completo. O \textbf{simulador} apenas reproduz o comportamento do sistema, executando o app compilado para o computador.
      \end{definicao}
      \begin{pordentro}[Por que o emulador exige virtualização]
        O emulador Android roda uma imagem completa do sistema numa máquina virtual. Com aceleração por hardware (Intel VT-x, AMD-V, KVM no Linux, Hypervisor no macOS), as instruções rodam direto no processador e o emulador fica rápido.
      \end{pordentro}
    \end{column}
  \end{columns}
\end{frame}

\begin{frame}{Depuração}
  \begin{columns}[T]
    \begin{column}{0.5\textwidth}
      \begin{definicao}[Depurador]
        Ferramenta que pausa o programa num ponto escolhido (\emph{breakpoint}) para inspecionar variáveis, a pilha de chamadas e executar passo a passo.
      \end{definicao}
      \scriptsize
      \begin{tabularx}{\linewidth}{@{}l L@{}}
        \toprule
        \cabecalho{Comando} & \cabecalho{O que faz} \\
        \midrule
        \emph{Step over} & executa a linha e para na próxima \\
        \emph{Step into} & entra na função chamada \\
        \emph{Step out} & termina a função atual e volta \\
        \emph{Resume} & continua até o próximo \emph{breakpoint} \\
        \bottomrule
      \end{tabularx}
    \end{column}
    \begin{column}{0.48\textwidth}
      \begin{conceito}[Método científico da depuração]
        \begin{enumerate}
          \item Reproduza o problema de forma confiável.
          \item Formule uma hipótese sobre a causa.
          \item Teste a hipótese (\emph{breakpoint}, \emph{log}, teste automatizado).
          \item Corrija e escreva um teste que impeça a volta do defeito.
        \end{enumerate}
      \end{conceito}
      \begin{dica}
        Leia a mensagem de erro e a pilha de chamadas inteiras: a primeira linha do \textbf{seu} código costuma apontar a causa.
      \end{dica}
    \end{column}
  \end{columns}
\end{frame}
